% BHU PYQ -- Manually transcribed & error-corrected
% Paper: BCF-222 : Financial Analysis
% Exam: B.Com. (Hons.) FMM Semester IV Examination, 2022-23

\documentclass[11pt,a4paper]{article}
\usepackage[a4paper,top=2.5cm,bottom=2.5cm,left=2.8cm,right=2.8cm,headheight=14pt]{geometry}
\usepackage{amsmath,amssymb}
\usepackage[T1]{fontenc}
\usepackage[utf8]{inputenc}
\usepackage{lmodern,microtype,fancyhdr,enumitem,hyperref,lastpage,parskip}

\hypersetup{
    colorlinks=true,
    urlcolor=black,
    linkcolor=black,
    pdftitle={BCF-222: Financial Analysis -- BHU B.Com. (Hons.) FMM Sem IV 2022-23}
}

\pagestyle{fancy}
\fancyhf{}
\renewcommand{\headrulewidth}{0.4pt}
\renewcommand{\footrulewidth}{0.4pt}
\fancyhead[L]{\small   \textit{Banaras Hindu University}}
\fancyhead[R]{\small   \textit{BCF-222: Financial Analysis}}
\fancyfoot[C]{\small Page  \thepage\ of \pageref{LastPage}}


\newlist{parts}{enumerate}{2}
\setlist[parts,1]{label=(\alph*),leftmargin=*,itemsep=5pt,topsep=3pt}
\setlist[parts,2]{label=(\roman*),leftmargin=*,itemsep=3pt,topsep=2pt}

\newcommand{\pts}[1]{\hfill   \textit{[#1]}}


\begin{document}


\begin{center}
  {\large\bfseries BANARAS HINDU UNIVERSITY}\\[2pt]
  {
\normalsize B.Com. (Hons.) FMM --- Semester IV Examination, 2022-23}\\[6pt]
  \rule{0.8\linewidth}{0.5pt}\\[6pt]
  {\large\bfseries Paper: BCF-222: Financial Analysis}\\[4pt]
  {
\normalsize Time: 3 Hours\hfill Maximum Marks: 70}
\end{center}

\rule{\linewidth}{0.4pt}

\smallskip



\noindent   \textit{Instructions: Answer all questions. All questions carry equal marks (14 marks each).}

\bigskip




\noindent   \textbf{Question 1.} \\
What do you understand by financial statement analysis? Describe the importance of such statement.\pts{14}

\centerline{   \textbf{OR}}

What is common-size balance sheet and income statement? Explain the technique of preparing the common size balance sheet.\pts{14}

\medskip\rule{\linewidth}{0.2pt}




\noindent   \textbf{Question 2.} \\
Find out:

\begin{parts}
  \item Current Assets
  \item Current Liabilities
  \item Liquid Assets
\end{parts}
   \textbf{Given Data:}

\begin{itemize}[leftmargin=*]
  \item Current ratio = $2.6$
  \item Acid test ratio = $1.5$
  \item Working Capital = Rs. 1,64,000
\end{itemize}\pts{14}

\centerline{   \textbf{OR}}

Elucidate the limitations of financial analysis.\pts{14}

\medskip\rule{\linewidth}{0.2pt}




\noindent   \textbf{Question 3.} \\
What is funds flow statement? Give specimen of schedule of changes in the working capital and funds flow statement.\pts{14}

\centerline{   \textbf{OR}}

From the following balance sheet of XYZ Company for the year ending 31st Dec., 2021 and 31st Dec., 2022, prepare schedule of change in working capital and a statement showing sources and application of funds:\pts{14}


\begin{center}
\renewcommand{\arraystretch}{1.2}
\small

\begin{tabular}{|l|r|r||l|r|r|}
\hline
   \textbf{Liabilities} &   \textbf{2021 (Rs.)} &   \textbf{2022 (Rs.)} &   \textbf{Assets} &   \textbf{2021 (Rs.)} &   \textbf{2022 (Rs.)} \\ \hline
Share Capital        & 3,00,000            & 4,00,000            & Machinery       & 50,000              & 60,000              \\
Creditors            & 1,00,000            & 70,000              & Furniture       & 10,000              & 15,000              \\
P/L A/c              & 15,000              & 30,000              & Stock           & 85,000              & 1,05,000            \\
                     &                     &                     & Debtors         & 1,60,000            & 1,50,000            \\
                     &                     &                     & Cash            & 1,10,000            & 1,70,000            \\ \hline
   \textbf{Total}       &   \textbf{4,15,000}   &   \textbf{5,00,000}   &   \textbf{Total}  &   \textbf{4,15,000}   &   \textbf{5,00,000}   \\ \hline
\end{tabular}
\end{center}

\medskip\rule{\linewidth}{0.2pt}




\noindent   \textbf{Question 4.} \\
What do you mean by 'Cash Flow Statement'? How is it different from 'Fund Flow Statement'? What is its utility?\pts{14}

\centerline{   \textbf{OR}}

What is the purpose of preparing a cash flow statement? How is it prepared? Explain with example.\pts{14}

\medskip\rule{\linewidth}{0.2pt}




\noindent   \textbf{Question 5.} \\
What is leverage? Explain the operating leverage and financial leverage. How are these calculated?\pts{14}

\centerline{   \textbf{OR}}

Discuss the importance of the following terms in relation to marginal costing (any three):\pts{14}

\begin{parts}
  \item P/V Ratio
  \item Angle of incidence
  \item Margin of safety
  \item BEP (Break-Even Point)
\end{parts}

\vfill
\rule{\linewidth}{0.4pt}

\begin{center}\small
  \href{https://bhu-exam.vercel.app/}{Click here to visit our website}\\[4pt]
  \href{https://me-aryan.vercel.app/}{Click here to visit our contributor}\\[4pt]
  \href{https://quashan-me.vercel.app/}{Click here to visit our contributor}
\end{center}

\end{document}
